\section{Adaptive machines}\label{sec:adaptive}

An adaptive machine has many traces and no common dependency DAG. Two routes lead from the oblivious results to adaptive machines.
The first uses a single heavy trace of the machine itself; the second passes through the bit simulator of the fifteenth paper, whose
traces are few.

\begin{theorem}[Theorem~E, heavy traces and simulators]\label{thm:E}
\formal{TightConstants.lean}{frozenTrace_crossingLowerGT2H_menger_lb_tight, frozenTrace_macroscopicCrossingLowers_menger_lb_tight, frozenTrace_reduces_to_crossingCapableSkeleton_and_singleLeg_tight, simulator_heavyTrace_expansion_and_skeleton_tight, adaptive_simulator_heavyTrace_macroscopicExpansion_tight}
\begin{enumerate}[label=(\alph*)]
\item \emph{Heavy traces.} Let $M$ be any machine with $k$ work tapes, $C\ge1$ and $\alpha>0$, put $\theta=\thetaM(M,C,\alpha)$ and
$\rho_0=\rho_0(k,C,\tfrac{19}{20}\alpha)$, and let $d\ge\max\bigl(1,d_{\mathrm{br}}(k,C,\tfrac{19}{20}\alpha)\bigr)$. Suppose $M$ transposes at size
$d$ within time $C$ (as it does if it transposes within a smaller time), and let $\tau$ be an $\eta$-heavy trace of $M$. Then for every
$1\le H\le\theta L$ and every set $V_0$ of at least $\alpha L$ input positions, every set separating $\Out(V_0)$ from $\CL_H$ in the
dependency DAG of $\tau$ has at least
\[
  \cE(M,C,\alpha,\eta)\,\abs{V_0},\qquad \cE(M,C,\alpha,\eta)=\frac{\rho_0/2-\eta/\alpha}{\wval},
\]
elements; at $\eta=0$ this is $\cE(M,C,\alpha)$. If moreover $\eta\le(\alpha/8)\rho_0$ and $d\ge d_0$ with $d_0$ as in \cref{thm:X}, the
conclusion of \cref{thm:X} at $H^*$ holds in the dependency DAG of $\tau$ with the expansion constant
$\cE'(M,C,\alpha)=\frac38\rho_0/\wval\le\cE(M,C,\alpha,\eta)$ in place of $\cE(M,C,\alpha)$; if in addition $\alpha\le1$ and
$d\ge\lceil(8/\theta)^{1/2}\rceil$, so do the conclusions (a) and (b) of \cref{thm:S}, with $\cE'$ in place of $\cE$ in (b).
\item \emph{Simulators.} Let $M$ be any machine with $k$ work tapes, $c>0$ and $0<\alpha\le1$, and put $\eta^*=(\alpha/16)\rho_0(k,11c+1,\tfrac{19}{20}\alpha)$.
There are $s\ge1$ and $d_{\mathrm{sim}}$ such that the simulator $\Ns$ of~\cite[Theorem~S]{HighRate2026}, which has the same work tapes and
work alphabet as $M$ and computes the same partial function, satisfies the following at every $d\ge d_{\mathrm{sim}}$ at which $M$
transposes within time $c$: $\Ns$ transposes within time $11c+1$; it has at most $2^{\eta^*L}$ distinct traces on the inputs of length
$L$; its $2\eta^*$-heavy traces are followed together by more than $(1-2^{-\eta^*L})2^L$ inputs; and on each of its $2\eta^*$-heavy traces,
with $\theta$ and $H^*$ computed for $\Ns$ and the time constant $11c+1$, $W\le H^*$, $\theta L\le2H^*$, the separation bound of \cref{thm:X} at
$H^*$, and the conclusions of \cref{thm:S}(a),(b) all hold with the expansion constant $\cE'$ of $\Ns$.
\end{enumerate}
\end{theorem}

\begin{outline}
(a) \Cref{lem:frozen} gives $\abs U-\wval\abs\Gamma\le\eta L$ for every set $\Gamma$ of steps and every set $U$ acyclic in the cone
relation of $\tau$ avoiding $\Gamma$. The proof of \cref{thm:X} uses the transposition hypothesis only through such a bound, with $0$ in
place of $\eta L$, and gives $\wval\abs B\ge\frac12\rho_0\abs{V_0}-\eta L\ge(\frac12\rho_0-\eta/\alpha)\abs{V_0}$, since
$\alpha L\le\abs{V_0}$ (\lean{linearTimeTrace_crossingLowerGT2H_menger_lb_of_coneDeficit_tight}). For $\eta\le(\alpha/8)\rho_0$ this is at
least $\frac38\rho_0\abs{V_0}$. (b) The constant $\eta^*$ depends only on $k$, $c$
and $\alpha$. The fifteenth paper's trace transfer at $\eta^*$ (\lean{TapeStorageRates.aligned_low_trace_dimensions}) gives $s$ and the
simulator, with at most $2^{\eta^*L}$ traces; the light traces are followed by fewer than $2^{\eta^*L}2^{(1-2\eta^*)L}$ inputs. Since $\Ns$ has
$k$ work tapes and runs in time $11c+1$, and $2\eta^*=(\alpha/8)\rho_0(k,11c+1,\tfrac{19}{20}\alpha)$, part (a) applies to $\Ns$ on each heavy
trace.
The quantifiers are ordered $\eta^*$, then $s$ and $\Ns$, then the constants of $\Ns$ and $d_{\mathrm{sim}}$.
\end{outline}

Part (b) is the only route by which the macroscopic condition reaches arbitrary linear-time transposers so far, and it reaches them
through a different machine's dependency DAG. The simulator rewrites windows of $3s$ cells in each round, and a separating set in its
DAG need not correspond to one in the DAG of $M$ (\cref{prob:pullback}). Its value width grows with $s$, so its constants are much smaller
than those of $M$. Statements in this paper about the dependency DAG of a general machine therefore always assume a heavy trace or name the
simulator.

\begin{remark}[The heaviness hypothesis cannot simply be dropped]\label{rem:rare}
Proof sketch only: given any machine $M_0$, add three states so that the new machine first scans the input for a $1$; on the all-zero
input it writes the all-zero output with its work heads standing still, and otherwise it returns its input head and runs $M_0$. The
new machine transposes whenever $M_0$ does, within a constant more time, and its all-zero trace makes no revisits at all. Hence the
version of \cref{thm:E}(a) without a lower bound on the trace probability would exclude this machine, and so would imply $\neg$(T).
It is not refuted: it is false if (T) holds.
\end{remark}

\begin{remark}[Wave counts on heavy traces]\label{rem:waves}
Because every expanding set of steps is as large as its expansion, $m$ pairwise disjoint sets of steps into which a set of $\alpha L$ outputs
expands by a factor $c$ need time at least $mc\alpha L$ (\lean{disjointWave_mengerExpansion_counting_le}); for oblivious machines in time
$C(L+1)$ there are at most $2C/(c\alpha)$ of them (\lean{obliviousLinTime_expandingWaves_le}). A hypothesis asserting unboundedly many such
sets on every heavy trace of every linear-time transposer is therefore equivalent to there being no linear-time transposer at
infinitely many sizes, and gives no route to $\neg$(T). For oblivious machines, a hypothesis asserting unboundedly many such sets on every
transposer, of any running time, would exclude oblivious linear-time transposers; the converse is not known.
\end{remark}

The fibre argument needs correctness only on the inputs it counts, so \cref{thm:X} and the separation bound of \cref{thm:E}(a)
extend to machines that transpose correctly at least a $2^{-\eta L}$ fraction of their inputs. The constant $\cE(M,C,\alpha,\eta)$
of \cref{thm:E}(a) equals $\cE(M,C,\alpha)-\eta/(\alpha\wval)$, and it is positive exactly for
$\eta<\frac12\alpha\rho_0(k,C,\tfrac{19}{20}\alpha)$.

\begin{corollary}[Approximate transposers]\label{cor:approx}
\formal{TightConstants.lean}{correctSubset_crossingLowerGT2H_menger_lb_tight, obliviousOn_correctSubset_macroscopicCrossingLowers_menger_lb_tight}
Let $M$ be any machine with $k$ work tapes, $C\ge1$, $\alpha>0$ and $\eta\ge0$.
\begin{enumerate}[label=(\alph*)]
\item Let $d\ge\max\bigl(1,d_{\mathrm{br}}(k,C,\tfrac{19}{20}\alpha)\bigr)$, and let $\tau$ be a trace of $M$ followed by at least
$2^{(1-\eta)L}$ inputs that $M$ transposes within time $C$. For every $1\le H\le\thetaM(M,C,\alpha)L$ and every set $V_0$ of at least
$\alpha L$ input positions, every set of steps separating $\Out(V_0)$ from $\CL_H$ in the dependency DAG of $\tau$ has at least
$\cE(M,C,\alpha,\eta)\abs{V_0}$ elements.
\item Let $d\ge d_0$ with $d_0$ as in \cref{thm:X}, and let $M$ be oblivious at length $L$ with halting time $T$ and transpose at least
$2^{(1-\eta)L}$ inputs within time $C$. Then the conclusion of \cref{thm:X} at $H^*$ holds in the common trace, with $\cE(M,C,\alpha,\eta)$
in place of $\cE(M,C,\alpha)$; for other $H$, (a) applies, since every input follows the common trace.
\end{enumerate}
\end{corollary}

\begin{outline}
(a) The proof of \cref{thm:E}(a) goes through with \cref{lem:frozen} applied to the correctly transposed inputs that follow $\tau$; the
deficit $\eta L\le(\eta/\alpha)\abs{V_0}$ lowers the constant by $\eta/(\alpha\wval)$. (b) All inputs follow the common trace, which halts
at $T$, and any correctly transposed input gives $T\le C(L+1)$; apply (a) at $H^*$. With $\eta=0$, (b) is \cref{thm:X}.
\end{outline}

Read the other way, a certificate on each of few traces bounds the fraction of inputs that a machine transposes correctly.

\begin{corollary}[Deficient traces carry few correct inputs]\label{cor:frozen-count}
\formal{SimulatorBridge.lean}{card_correctInputs_le_of_coneDeficit}
Let $M$ be any machine, $d\ge1$, $c>0$, and $\mathcal T$ a finite set of traces of $M$, each with a \emph{deficit certificate}: a set
$\Gamma_\tau$ of steps and a set $U_\tau$ acyclic in the cone relation of $\tau$ avoiding $\Gamma_\tau$ with
$\abs{U_\tau}-\wval\abs{\Gamma_\tau}\ge\delta L$. Then at most $\abs{\mathcal T}\,2^{(1-\delta)L}$ inputs follow a trace in $\mathcal T$
and are transposed by $M$ within time $c$. If $\abs{\mathcal T}\le2^{\eta L}$, these inputs are at most a $2^{-(\delta-\eta)L}$ fraction of
all inputs.
\end{corollary}

\begin{outline}
For each $\tau\in\mathcal T$ followed by correctly transposed inputs, \cref{lem:frozen} applied to the set $S_\tau$ of these inputs gives
$\delta L\le L-\log_2\abs{S_\tau}$. Sum $\abs{S_\tau}\le2^{(1-\delta)L}$ over $\mathcal T$.
\end{outline}
